% Options for packages loaded elsewhere
\PassOptionsToPackage{unicode}{hyperref}
\PassOptionsToPackage{hyphens}{url}
\PassOptionsToPackage{dvipsnames,svgnames,x11names}{xcolor}
\PassOptionsToPackage{space}{xeCJK}
\documentclass[
]{article}
\usepackage{xcolor}
\usepackage[margin=2.5cm]{geometry}
\usepackage{amsmath,amssymb}
\setcounter{secnumdepth}{-\maxdimen} % remove section numbering
\usepackage{iftex}
\ifPDFTeX
  \usepackage[T1]{fontenc}
  \usepackage[utf8]{inputenc}
  \usepackage{textcomp} % provide euro and other symbols
\else % if luatex or xetex
  \usepackage{unicode-math} % this also loads fontspec
  \defaultfontfeatures{Scale=MatchLowercase}
  \defaultfontfeatures[\rmfamily]{Ligatures=TeX,Scale=1}
\fi
\usepackage{lmodern}
\ifPDFTeX\else
  % xetex/luatex font selection
  \setmainfont[]{DejaVu Sans}
  \setmonofont[]{DejaVu Sans Mono}
  \ifXeTeX
    \usepackage{xeCJK}
    \setCJKmainfont[]{Noto Sans CJK SC}
  \fi
  \ifLuaTeX
    \usepackage[]{luatexja-fontspec}
    \setmainjfont[]{Noto Sans CJK SC}
  \fi
\fi
% Use upquote if available, for straight quotes in verbatim environments
\IfFileExists{upquote.sty}{\usepackage{upquote}}{}
\IfFileExists{microtype.sty}{% use microtype if available
  \usepackage[]{microtype}
  \UseMicrotypeSet[protrusion]{basicmath} % disable protrusion for tt fonts
}{}
\makeatletter
\@ifundefined{KOMAClassName}{% if non-KOMA class
  \IfFileExists{parskip.sty}{%
    \usepackage{parskip}
  }{% else
    \setlength{\parindent}{0pt}
    \setlength{\parskip}{6pt plus 2pt minus 1pt}}
}{% if KOMA class
  \KOMAoptions{parskip=half}}
\makeatother
\usepackage{longtable,booktabs,array}
\usepackage{caption}
\captionsetup[table]{skip=6pt}
\newcounter{none} % for unnumbered tables
\usepackage{calc} % for calculating minipage widths
% Correct order of tables after \paragraph or \subparagraph
\usepackage{etoolbox}
\makeatletter
\patchcmd\longtable{\par}{\if@noskipsec\mbox{}\fi\par}{}{}
\makeatother
% Allow footnotes in longtable head/foot
\IfFileExists{footnotehyper.sty}{\usepackage{footnotehyper}}{\usepackage{footnote}}
\makesavenoteenv{longtable}
\setlength{\emergencystretch}{3em} % prevent overfull lines
\providecommand{\tightlist}{%
  \setlength{\itemsep}{0pt}\setlength{\parskip}{0pt}}
\usepackage{bookmark}
\IfFileExists{xurl.sty}{\usepackage{xurl}}{} % add URL line breaks if available
\urlstyle{same}
% fallback for those not using the hyperref driver hyperxmp:
\makeatletter
\@ifundefined{xmpquote}{\newcommand{\xmpquote}[1]{#1}}{}
\makeatother
\hypersetup{
  colorlinks=true,
  linkcolor={Maroon},
  filecolor={Maroon},
  citecolor={Blue},
  urlcolor={Blue},
  pdfcreator={LaTeX via pandoc}}

\author{}
\date{}

\begin{document}

\section{筑基篇课后习题 VI：定义 pat
的触发------剥离一层自指（换基点观测实验）}\label{ux7b51ux57faux7bc7ux8bfeux540eux4e60ux9898-viux5b9aux4e49-pat-ux7684ux89e6ux53d1ux5265ux79bbux4e00ux5c42ux81eaux6307ux6362ux57faux70b9ux89c2ux6d4bux5b9eux9a8c}

\begin{quote}
\textbf{声明 (Declaration)}: 本文\textbf{不代表任何数学结论}
(non-mathematical-claim), 仅为基于 Lean 4
形式化的一种\textbf{实验观测报告} (experimental observation report)。

{\def\LTcaptype{none} % do not increment counter
\begin{longtable}[]{@{}
  >{\raggedright\arraybackslash}p{(\linewidth - 6\tabcolsep) * \real{0.2500}}
  >{\raggedright\arraybackslash}p{(\linewidth - 6\tabcolsep) * \real{0.2500}}
  >{\raggedright\arraybackslash}p{(\linewidth - 6\tabcolsep) * \real{0.2500}}
  >{\raggedright\arraybackslash}p{(\linewidth - 6\tabcolsep) * \real{0.2500}}@{}}
\toprule\noalign{}
\begin{minipage}[b]{\linewidth}\raggedright
习题
\end{minipage} & \begin{minipage}[b]{\linewidth}\raggedright
主题
\end{minipage} & \begin{minipage}[b]{\linewidth}\raggedright
Lean 形式化
\end{minipage} & \begin{minipage}[b]{\linewidth}\raggedright
DOI
\end{minipage} \\
\midrule\noalign{}
\endhead
\bottomrule\noalign{}
\endlastfoot
exercise\_i\_4 & 触发-剥离实验（自指坍缩的启动机制） &
PatTriggerExperiment.lean & 10.5281/zenodo.21916845 \\
\end{longtable}
}
\end{quote}

\textbf{Exercise VI for the Foundation-Building Chapter: Defining the
Trigger of PAT --- Peeling One Layer of Self-Reference
(Basepoint-Changed Observation Experiment)}

\begin{quote}
习题定位：筑基篇课后习题系列第六题（承接习题 V
杨-米尔斯质量间隙的语境------一直在用 pat 但未定义 pat
的触发）。用户假设（2026-08-13）：\textbf{pat
的触发本质上可能是剥离一层自指}。本习题把该假设精确化、设计实验（换基点观测剥离过程）、Lean
静态定理验收 + 数值模拟观测。诚实边界：触发 =
剥离是结构假设（对合族内验证），不回答''方向声明从何而来''的终极来源。
\end{quote}

\emph{2026-08-13 · Internal research exercise · Lean 4 / mathlib v4.32.2
· 7 new theorems PROVED no sorry + 模拟 4/4 PASS ·
算器神魂论筑基篇课后习题 VI }

\begin{center}\rule{0.5\linewidth}{0.5pt}\end{center}

\subsection{习题陈述}\label{ux4e60ux9898ux9648ux8ff0}

\begin{quote}
定义 pat 的触发。筑基篇已给出：方向声明（R136）= 链的启动条件，未声明 =
自指坍缩（R122/R134）；但''声明''从何而来？用户假设：\textbf{触发 =
剥离一层自指}------自指环路（对合，无净移动）经剥离提出相位层（位移场，方向显现），方向声明
⟹ 链展开。实验：换基点观测剥离自指的过程，验证三预言（P1 相位显现 / P2
自相似 / P3 换基点不变）。
\end{quote}

\begin{center}\rule{0.5\linewidth}{0.5pt}\end{center}

\subsection{零、总论：为什么需要定义 pat
的触发}\label{ux96f6ux603bux8bbaux4e3aux4ec0ux4e48ux9700ux8981ux5b9aux4e49-pat-ux7684ux89e6ux53d1}

筑基篇用 pat 构造了一切（数域/素数/连续统/计算），但 pat
的启动条件只停在''方向声明''（R136）------声明本身是一个操作，它的机制是什么？已有线索：

\begin{itemize}
\tightlist
\item
  \textbf{未声明 = 坍缩}（R122）：自指 IS pat n------自指环路（选择方向
  = 互逆方向）相位无净移动 ⟹ 坍缩到 pat 0。R134：pat0
  吸收（防自指定义污染）。
\item
  \textbf{拆自指内部还是
  SRT}（R129）：每拆一层出现另一层（自指自相似）。
\item
  \textbf{剥离方向单步不可判定}（R130）：深入（Δ→0）/平移（Δ恒定）/远离（Δ增大），单次剥离无法区分。
\end{itemize}

用户的合成：\textbf{触发（从坍缩进入链）=
剥离一层自指}------剥离显现相位差（方向，RulerPhase），方向声明（R136）成为剥离的必然产物，而非外来的操作。

\begin{center}\rule{0.5\linewidth}{0.5pt}\end{center}

\subsection{一、定义：自指结构 + 剥离 +
触发（精确化）}\label{ux4e00ux5b9aux4e49ux81eaux6307ux7ed3ux6784-ux5265ux79bb-ux89e6ux53d1ux7cbeux786eux5316}

全部实数（命名纪律：无 sqrt、无无声明的 i）：

\begin{itemize}
\tightlist
\item
  \textbf{自指结构} = 基点镜像
  \texttt{S\_e(x)\ =\ 2e\ -\ x}（对合：S\_e(S\_e x) = x，R128
  任意基点镜像对合；R085 折叠类）。自指环路 = 对合迭代，无净移动。
\item
  \textbf{剥离（peel）} = 从对合提取位移场
  \texttt{Δ(x)\ =\ S\_e(x)\ -\ x\ =\ 2(e-x)}------相位差（方向，RulerPhase；R130：剥离是相位操作）。
\item
  \textbf{触发} = 剥离：方向显现（Δ ≠ 0 当 x ≠ e）⟹ 方向声明（R136）⟹
  链展开（R137）。未剥离（对合黑箱）= 无净移动 = 坍缩到 pat
  0（R122/R134）。
\end{itemize}

\subsubsection{形式化（PatTriggerExperiment.lean，7
定理）}\label{ux5f62ux5f0fux5316pattriggerexperiment.lean7-ux5b9aux7406}

\begin{itemize}
\tightlist
\item
  \texttt{selfref\_no\_net\_motion}：S\_e(S\_e x) =
  x------自指环路无净移动（未剥离 = 坍缩）。
\item
  \texttt{peel\_reveals\_phase}：S\_e x - x =
  2(e-x)------剥离显现位移场。
\item
  \texttt{peel\_phase\_nonzero\_off\_basepoint}：x ≠ e ⟹ Δ ≠ 0------触发
  = 剥离 ⟹ 方向显现（P1）。
\item
  \texttt{peel\_displacement\_flips\_under\_self}：Δ(S x) =
  -Δ(x)------剥离的剥离 = 相位取反（P2 自相似，R129/R074）。
\item
  \texttt{peel\_basepoint\_translation}：S\_f(x+(f-e)) =
  S\_e(x)+(f-e)------换基点 = 平移共轭（P3，R128）。
\item
  \texttt{peel\_displacement\_basepoint\_invariance}：Δ\_f(x+(f-e)) =
  Δ\_e(x)------位移场换基点平移不变（P3 核心）。
\item
  \texttt{trigger\_peel\_perspective}：全景组合。
\end{itemize}

\textbf{验证}：lake env lean 编译通过，0 error 0 sorry；恒等式方向
heap\_verify 枚举 5/5 PASS（Z/7：对合/位移场/平移共轭/换基点不变）。

\begin{center}\rule{0.5\linewidth}{0.5pt}\end{center}

\subsection{二、三预言（可证伪）}\label{ux4e8cux4e09ux9884ux8a00ux53efux8bc1ux4f2a}

{\def\LTcaptype{none} % do not increment counter
\begin{longtable}[]{@{}
  >{\raggedright\arraybackslash}p{(\linewidth - 6\tabcolsep) * \real{0.2500}}
  >{\raggedright\arraybackslash}p{(\linewidth - 6\tabcolsep) * \real{0.2500}}
  >{\raggedright\arraybackslash}p{(\linewidth - 6\tabcolsep) * \real{0.2500}}
  >{\raggedright\arraybackslash}p{(\linewidth - 6\tabcolsep) * \real{0.2500}}@{}}
\toprule\noalign{}
\begin{minipage}[b]{\linewidth}\raggedright
预言
\end{minipage} & \begin{minipage}[b]{\linewidth}\raggedright
内容
\end{minipage} & \begin{minipage}[b]{\linewidth}\raggedright
锚到
\end{minipage} & \begin{minipage}[b]{\linewidth}\raggedright
结果
\end{minipage} \\
\midrule\noalign{}
\endhead
\bottomrule\noalign{}
\endlastfoot
\textbf{P1 相位显现} & 剥离 ⟹ Δ(x) ≠ 0（x ≠
e）；未剥离（对合环路）无净移动 & R122/R134 + RulerPhase & Lean + 模拟
PASS \\
\textbf{P2 自相似} & Δ(S\_e x) = -Δ(x)------剥离的剥离 =
相位取反，内部还是一套 SRT & R129 + R074 & Lean + 模拟 PASS \\
\textbf{P3 换基点不变} & Δ\_f(x+(f-e)) =
Δ\_e(x)------剥离结果与基点选择无关（只差平移） & R128 平移共轭 & Lean +
模拟 PASS \\
\end{longtable}
}

P3
是''换基点观测''的核心：\textbf{剥离自指的过程在换基点下平移不变}------在基点
0、1、3、-2 剥离同一个 x，位移场只是平移（数值模拟见表）。

\begin{center}\rule{0.5\linewidth}{0.5pt}\end{center}

\subsection{三、换基点观测表（模拟，R161\_trigger\_peel.py，2000 样本 ×
4
预言）}\label{ux4e09ux6362ux57faux70b9ux89c2ux6d4bux8868ux6a21ux62dfr161_trigger_peel.py2000-ux6837ux672c-4-ux9884ux8a00}

\begin{verbatim}
P1 触发=剥离显现方向: Δ≠0 当 x≠e: PASS
P1b 未剥离 (对合环路) 无净移动: PASS
P2 自相似: Δ(S x) = -Δ(x) (剥离的剥离=相位取反): PASS
P3 换基点平移不变: Δ_f(x+(f-e)) = Δ_e(x): PASS

换基点观测表 (剥离 S_e 于同一 x=1.5):
  基点 e=0.0: Δ = -3.0000 (方向: 指向基点)
  基点 e=1.0: Δ = -1.0000 (方向: 指向基点)
  基点 e=3.0: Δ = +3.0000 (方向: 指向基点)
  基点 e=-2.0: Δ = -7.0000 (方向: 指向基点)
\end{verbatim}

\textbf{观测结论}：剥离的位移场\textbf{总是指向基点}------换基点不改变剥离的语义（方向指向基点），只改变位移的大小（\textbar Δ\textbar{}
= 2\textbar x-e\textbar，距基点 2 倍）。这正是 R128''原点现象 ≡
基点现象（平移共轭）``在剥离语境下的实例：\textbf{触发机制是基点无关的}。

\begin{center}\rule{0.5\linewidth}{0.5pt}\end{center}

\subsection{四、趋势观测（R130
的序列版）}\label{ux56dbux8d8bux52bfux89c2ux6d4br130-ux7684ux5e8fux5217ux7248}

\begin{verbatim}
趋势观测: x 从基点出发渐远 (x = e + k·0.5), |Δ| 行为:
  镜像 S_e (自指): |Δ| = [1.00, 2.00, 3.00, 4.00, 5.00] → 远离 (|Δ|增大)
  反演 I (自指):   |Δ| = [1.00, 2.00, 3.00, 4.00, 5.00] → 远离 (|Δ|增大)
  平移 T_d (对照): |Δ| = [1.00, 1.00, 1.00, 1.00, 1.00] → 平移 (|Δ|恒定)
\end{verbatim}

\textbf{观测结论}： - 自指结构（对合）在 x 序列上剥离的
\textbar Δ\textbar{} = 2\textbar x-e\textbar{}
\textbf{线性增大}（远离趋势）； -
平移（对照组，非自指）\textbar Δ\textbar{} 恒定（平移趋势）； -
\textbf{观测方向决定趋势}：远离基点看 = 远离；逼近基点看（x →
e）\textbar Δ\textbar{} → 0 = 深入。R130
的单步不可判定性保持（单步剥离仍不可判定），但\textbf{在 x
序列上趋势可判定，且趋势是观测方向的选择}------这正是
R130''需要完整相位序列''的序列版。

\begin{center}\rule{0.5\linewidth}{0.5pt}\end{center}

\subsection{五、全景：触发 =
剥离一层自指}\label{ux4e94ux5168ux666fux89e6ux53d1-ux5265ux79bbux4e00ux5c42ux81eaux6307}

\textbf{★核心：未剥离（对合黑箱）= 自指环路无净移动 ⟹ 坍缩到 pat
0（R122/R134）；剥离 = 位移场显现方向（P1）⟹ 方向声明（R136）⟹
链展开（R137）；剥离的剥离 = 相位取反（P2 自相似，R129------内部还是一套
SRT）；换基点 =
平移共轭，剥离结果平移不变（P3，R128------触发机制基点无关）。}

\begin{itemize}
\tightlist
\item
  \texttt{trigger\_peel\_perspective}（Lean）：① 对合无净移动 ∧ ②
  剥离显现方向非零 ∧ ③ 位移场自指取反 ∧ ④ 换基点平移不变。
\end{itemize}

\begin{center}\rule{0.5\linewidth}{0.5pt}\end{center}

\subsection{习题解答总结}\label{ux4e60ux9898ux89e3ux7b54ux603bux7ed3}

{\def\LTcaptype{none} % do not increment counter
\begin{longtable}[]{@{}
  >{\raggedright\arraybackslash}p{(\linewidth - 6\tabcolsep) * \real{0.2500}}
  >{\raggedright\arraybackslash}p{(\linewidth - 6\tabcolsep) * \real{0.2500}}
  >{\raggedright\arraybackslash}p{(\linewidth - 6\tabcolsep) * \real{0.2500}}
  >{\raggedright\arraybackslash}p{(\linewidth - 6\tabcolsep) * \real{0.2500}}@{}}
\toprule\noalign{}
\begin{minipage}[b]{\linewidth}\raggedright
论点
\end{minipage} & \begin{minipage}[b]{\linewidth}\raggedright
内容
\end{minipage} & \begin{minipage}[b]{\linewidth}\raggedright
锚到
\end{minipage} & \begin{minipage}[b]{\linewidth}\raggedright
证据
\end{minipage} \\
\midrule\noalign{}
\endhead
\bottomrule\noalign{}
\endlastfoot
未剥离 = 坍缩 & 对合环路无净移动（自指 IS pat n） & R122/R134 &
selfref\_no\_net\_motion + P1b PASS \\
触发 = 剥离显现方向 & Δ = 2(e-x) ≠ 0（x ≠ e） & RulerPhase + R136 &
peel\_reveals\_phase / peel\_phase\_nonzero + P1 PASS \\
剥离自相似 & Δ(S x) = -Δ(x)（内部还是 SRT） & R129/R074 &
peel\_displacement\_flips\_under\_self + P2 PASS \\
换基点不变 & 剥离结果平移共轭，总指向基点 & R128 &
peel\_basepoint\_translation / peel\_displacement\_invariance + P3
PASS \\
★趋势 & 对合序列远离（\textbar Δ\textbar{} =
2\textbar x-e\textbar），观测方向决定趋势 & R130 序列版 &
模拟趋势观测表 \\
\end{longtable}
}

\textbf{习题的回答（一句话）}：\textbf{pat 的触发 =
剥离一层自指------未剥离的自指环路（对合）无净移动坍缩到 pat 0；剥离一层
= 位移场显现（方向 =
相位差），方向声明成为剥离的必然产物，链由此展开；剥离过程换基点平移不变（总指向基点），剥离的剥离
= 相位取反（自指自相似，内部还是一套 SRT）。}

\textbf{诚实边界}：触发 = 剥离是结构假设（对合族内验证 7 定理 + 模拟 4/4
PASS），不回答''方向声明从何而来''的终极来源（R136 的机制假设）；R130
单步不可判定性保持（趋势判定需要 x 序列，且观测方向本身是选择）。

\begin{center}\rule{0.5\linewidth}{0.5pt}\end{center}

\subsection{六、相关对照}\label{ux516dux76f8ux5173ux5bf9ux7167}

\begin{quote}
完整书目见习题 I
\texttt{../shared/pat\_pnp\_shared\_references\_ZH.md}（共享引用清单）（全字段
+ 验证记录）。本条习题全部锚本框架：
\end{quote}

{\def\LTcaptype{none} % do not increment counter
\begin{longtable}[]{@{}
  >{\raggedright\arraybackslash}p{(\linewidth - 4\tabcolsep) * \real{0.3333}}
  >{\raggedright\arraybackslash}p{(\linewidth - 4\tabcolsep) * \real{0.3333}}
  >{\raggedright\arraybackslash}p{(\linewidth - 4\tabcolsep) * \real{0.3333}}@{}}
\toprule\noalign{}
\begin{minipage}[b]{\linewidth}\raggedright
框架 claim
\end{minipage} & \begin{minipage}[b]{\linewidth}\raggedright
对应内容
\end{minipage} & \begin{minipage}[b]{\linewidth}\raggedright
备注
\end{minipage} \\
\midrule\noalign{}
\endhead
\bottomrule\noalign{}
\endlastfoot
R122 & 自指 IS pat n；自指环路无净移动 ⟹ 坍缩 & 未剥离 = 坍缩（P1b） \\
R128 & 原点现象 ≡ 基点现象；S\_e = T∘S₀∘T⁻¹ 平移共轭 &
换基点观测（P3） \\
R129 & 拆自指内部还是一套 SRT（自指自相似） & 剥离的剥离 =
相位取反（P2） \\
R130 & SRT 剥离方向单步不可判定（深入/平移/远离） &
趋势观测（序列版） \\
R134 & pat0 吸收（防自指定义污染） & 未触发 = 吸收态 \\
R136 & 方向声明 = 链启动 & 触发 = 剥离 ⟹ 声明 \\
R137 & Pat N 构造（全向 0 出发，声明方向） & 触发后的链展开 \\
R085/R074 & 折叠类 0 / 自指迭代的自指迭代 & 基点处坍缩 / 方向对偶 \\
\end{longtable}
}

\begin{center}\rule{0.5\linewidth}{0.5pt}\end{center}

\emph{筑基篇课后习题 VI · 2026-08-13 · Lean 4 / mathlib v4.32.2 · 7
theorems PROVED no sorry（PatTriggerExperiment.lean）+ 模拟 4/4
PASS（R161\_trigger\_peel.py）· 配套 Lean
形式化：formal/Formal/Toolkit/PatTriggerExperiment.lean（快照见
../lean/PatTriggerExperiment.lean）}

\end{document}
